% PDF 127、130--133；原书第 97、100--103 页。
\section{提示工程}
上一节用几个简单的例子说明，如何利用提示解决不同的自然语言处理任务。提示的具体措辞对取得良好结果很重要。为提高模型在目标下游任务上的表现而设计合适的提示函数，就是提示工程。设计提示时，需要考虑并综合多个要素，包括选择预训练模型、确定合适的提示形式、设计提示模板和答案工程。模板的设计大致分为两类。
\begin{itemize}
\item 手工设计模板。
\item 自动设计模板。
\end{itemize}
前者借助人的专业知识和反复尝试来优化提示，后者用各种自动化方法寻找适合给定任务的模板。图~\ref{fig:ch03-design} 概览了接下来两个节的结构。下面介绍提示分类的基本术语，概述手工提示工程，并详细讨论文献中的几种自动化方法。

\begin{figure}[!htbp]
\centering
% 重绘自原书第 100 页的图 3.10。
\begin{tikzpicture}[
  font=\footnotesize\sffamily,
  option/.style={align=center,minimum width=24mm,minimum height=8mm,inner sep=2pt},
  wire/.style={draw=black!48,line width=.65pt}]
\node[option,fill=green!17] (root) at (0,0) {基于提示\\的学习};
\node[option,fill=yellow!65!orange] (model) at (2.92,3.80) {预训练模型};
\node[option,fill=blue!15] (eng) at (2.92,2.70) {提示工程};
\node[option,fill=orange!20] (ans) at (2.92,-2.00) {答案工程};
\node[option,fill=black!10] (multi) at (2.92,-3.10) {多提示学习};
\node[option,fill=black!55,text=white] (tuning) at (2.92,-4.20) {基于提示的微调};

\node[option,fill=blue!26] (promptshape) at (5.84,3.80) {提示形式};
\node[option,fill=blue!26] (template) at (5.84,1.60) {提示模板设计};
\node[option,fill=orange!30] (answershape) at (5.84,-.60) {答案形式};
\node[option,fill=orange!30] (search) at (5.84,-3.30) {答案搜索};

\node[option,fill=blue!37] (cloze) at (8.76,4.35) {完形填空提示};
\node[option,fill=blue!37] (prefix) at (8.76,3.25) {前缀提示};
\node[option,fill=blue!59!black,text=white] (manual) at (8.76,2.30) {手工设计};
\node[option,fill=blue!59!black,text=white] (auto) at (8.76,1.38) {自动设计};
\node[option,fill=orange!59] (token) at (8.76,.50) {词元};
\node[option,fill=orange!59] (span) at (8.76,-.60) {片段};
\node[option,fill=orange!59] (sentence) at (8.76,-1.70) {句子};
\node[option,fill=orange!68!black,text=white] (human) at (8.76,-2.75) {人工设计};
\node[option,fill=orange!68!black,text=white] (automated) at (8.76,-3.85) {自动答案搜索};

\draw[wire] (root.east) -- (1.46,0);
\draw[wire] (1.46,-4.20) -- (1.46,3.80);
\foreach \target/\y in {model/3.80,eng/2.70,ans/-2.00,multi/-3.10,tuning/-4.20} {
  \draw[wire] (1.46,\y) -- (\target.west);
}
\draw[wire] (eng.east) -- (4.38,2.70) |- (promptshape.west);
\draw[wire] (4.38,2.70) |- (template.west);
\draw[wire] (ans.east) -- (4.38,-2.00) |- (answershape.west);
\draw[wire] (4.38,-2.00) |- (search.west);

\draw[wire] (promptshape.east) -- (7.30,3.80) |- (cloze.west);
\draw[wire] (7.30,3.80) |- (prefix.west);
\draw[wire] (template.east) -- (7.30,1.60) |- (manual.west);
\draw[wire] (7.30,1.60) |- (auto.west);
\draw[wire] (answershape.east) -- (span.west);
\draw[wire] (7.30,-.60) |- (token.west);
\draw[wire] (7.30,-.60) |- (sentence.west);
\draw[wire] (search.east) -- (7.30,-3.30) |- (human.west);
\draw[wire] (7.30,-3.30) |- (automated.west);
\end{tikzpicture}

\caption{基于提示的学习中，各种设计选择及决策之间的关系。}
\label{fig:ch03-design}
\end{figure}

\subsection{提示形式}
提示模板大致有两种形式，即前缀提示（prefix prompt）和完形填空提示（cloze prompt）。下文用“提示形式”（prompt shape）统称这种区分。

\subsubsection{前缀提示}
在前缀提示中，输入及描述输入的文本都位于答案空位之前，形成前缀。例如，在影评情感分析任务中，前缀提示模板可以写作
\begin{equation}
\text{“}\symbf{x}'=\text{[x] This movie is [z]}\text{”},
\tag{3.5}
\end{equation}
将模板应用于输入，会得到“Cannot watch this movie. This movie is [z]”这样的提示。

如果回答与输入可以按从左到右的顺序直接对应，两者之间的联系简单明确，前缀提示往往表现良好。例如，在机器翻译任务中，明确给出目标语言，前缀提示就可以有效地引导模型生成译文。考虑以下模板。
\begin{equation}
\begin{aligned}
\text{“}\symbf{x}'={}&\text{Translate the following English sentence}\\
 &\text{to French: [x][z]}\text{”},
\end{aligned}
\tag{3.6}
\end{equation}
前缀提示通过明确指定所需语言，向模型给出直接的指令，从而产生准确的译文。

\subsubsection{完形填空提示}
与前缀提示不同，完形填空提示的答案空位两侧都有模板词元，所以答案位于模板内部。例如，在影评情感分析任务中，完形填空提示模板可以表示为
\begin{equation}
\text{“}\symbf{x}'=\text{[x] This is a [z] movie.}\text{”},
\tag{3.7}
\end{equation}
答案空位由上下文中的词语包围。把模板应用到输入上，就会得到提示“Cannot watch this movie. This is a [z] movie.”。这里的标点也算模板文本。因此，若模板在答案空位 \texttt{[z]} 之后只有一个句号，仍视为完形填空提示。

完形填空提示形式灵活，可以构造出更自然的语言结构，尤其适合多项选择问答等任务。把答案空位放在一个有上下文的短语中，可以促使模型生成符合周围语言模式的输出。答案两侧的上下文让模型能考虑输入中的语言特征和情境线索，最终生成更准确、更有意义的内容。

两种提示形式的区别如下。
\begin{itemize}
\item \textbf{前缀提示。}输入和模板文字都位于答案空位之前。例如“Cannot watch this movie. This movie is [z]”。
\item \textbf{完形填空提示。}输入和模板文字围绕着答案空位。例如“Cannot watch this movie. It was a [z] movie.”。
\end{itemize}

\subsection{手工设计模板}
现在转向模板的设计和优化。最直接的方法是手工设计模板，利用人的专业知识、直觉和对任务的理解，设计合适的模板。通常既要考虑任务本身的特点，也要用启发式方法确定合适的提示结构和措辞。为了改善表现，可能需要反复调整模板（\refcite{ch03-brown}{Brown et al., 2020}；\refcite{ch03-petroni}{Petroni et al., 2019}；\refcite{ch03-schick-2020a}{Schick and Schütze, 2020a},b）。最终应根据在有标签数据集上的表现作决定。

手工生成提示没有通用做法，最重要的原则是试验多个候选。以返回各国首都的提示为例，可以看到反复尝试的重要性。下面给出四个候选模板。
\begin{lstlisting}[caption={国家首都查询的提示模板},label={lst:ch03-capitals}]
1. "the capital city of [x] is [z] ."
2. "[z] is the capital city of [x] ."
3. "what is the capital city of [x]? It is [z] ."
4. "[z] is located in [x], and is its capital city ."
\end{lstlisting}

这四个模板看起来都说得通，效果是否一样？原书用 AllenNLP 掩码语言建模演示工具\footnote{\url{https://demo.allennlp.org/masked-lm}} 测试了一个输入样例。在这个演示工具中输入一个带掩码词元的句子，模型会返回最可能填入空位的若干词元。以 Poland 为输入 \texttt{[x]}，对每个模板预测前三个词元，结果见表~\ref{tab:ch03-capitals}。

\begin{table}[!htbp]
\centering
\caption{在 AllenNLP 掩码语言模型演示中使用清单~\ref{lst:ch03-capitals} 各模板的预测分数。给定 \texttt{[x]} 为“Poland”，每行给出填入 \texttt{[z]} 的前三个预测词元的概率，以百分数表示。}
\label{tab:ch03-capitals}
\small
\begin{tabular}{@{}lrrrrrr@{}}
\toprule
输入 & Warsaw & Kraków & Poznan & Poland & here & It\\
\midrule
模板 1 & 37.7 & 23.2 & 19.3 & — & — & —\\
模板 2 & 55.7 & 14.8 & 15.0 & — & — & —\\
模板 3 & 10.1 & — & — & 21.2 & 3.6 & —\\
模板 4 & 29.3 & — & 12.3 & — & — & 33.0\\
\bottomrule
\end{tabular}
\end{table}

模板 1 和模板 2 的首选词元是正确答案“Warsaw”，模板 2 对“Warsaw”的预测优势更大。值得注意的是，它们是四个模板中最简单直接的两个，都没有多个句子或累赘的从句安排。模板 3 的首选答案是“Poland”，模板 4 则预测代词“It”。两者的第二选择都是“Warsaw”，但显然没有像前两种模板那样有效地激发大语言模型中隐含的知识。

在手工设计提示时，设计者应当用多个模板和许多类似的有标签样例进行试验，以便作统计优化。应根据某项指标确定最优提示，例如采用选出首选预测的办法。
\begin{equation}
A(t_r,i)=
\frac{\sum_{\langle x,y\rangle\in\mathcal{R}}
 \delta\bigl(y=\arg\max_{y'}P_{\mathrm{LM}}(y'\mid x,t_{r,i})\bigr)}
{\mathcal{R}}.
\tag{3.8}
\end{equation}
这里 $\mathcal{R}$ 是由关系 $r$ 对应的带标签主语—宾语对组成的测试集。$\delta(\cdot)$ 是克罗内克的 $\delta$ 函数。当 $y$ 等于语言模型的首选预测时，函数返回 1，否则返回 0。最终选择在这组主语—宾语对训练样例上准确率最高的提示。
\transnote{原书式（3.8）的分母印为 $\mathcal{R}$。若计算准确率，分母是集合大小 $|\mathcal{R}|$。本段也保留了原书先称“测试集”、后称“训练样例”的措辞差异。}
